\chapter{Codimension-one points are \texorpdfstring{$\bar k$}{kbar}-rational}
\label{chap:RiemannRoch_ResidueFieldKbar}

% archon:covers AlgebraicJacobian/RiemannRoch/ResidueFieldKbar.lean

% Chapter-local fallbacks for the order-theoretic dimension macros (not yet in
% blueprint/src/macros/common.tex).  See the writer report "needs macros" note.
\providecommand{\Order}{\mathrm{Order}}
\providecommand{\krullDim}{\dim}
\providecommand{\coheight}{\operatorname{coheight}}
\providecommand{\height}{\operatorname{ht}}

This chapter collects the project-local commutative-algebra and topological
substrate behind the codimension-one rationality statement
\cref{lem:codimOne_point_residueField_eq_kbar}, proved at the end of this
chapter and consumed by the regularity substrate of
\cref{chap:RiemannRoch_OcOfD}.  At a closed point of an integral finite-type
curve over an algebraically closed field $\bar k$ the residue field is $\bar k$
itself; the file isolates the two ingredients --- the affine-local residue-field
computation (Zariski/Nullstellensatz plus algebraic-closure triviality) and the
topological ``coheight one $\Rightarrow$ closed point'' implication --- as
standalone supplementary lemmas, then glues them into the scheme-level statement.

\begin{lemma}[Finite-type field extension of $\bar k$ is trivial]
  \label{lem:algebraMap_bijective_of_finiteType}
  \lean{IsAlgClosed.algebraMap_bijective_of_finiteType}
  \leanok
  Let $\bar k$ be algebraically closed and $K$ a field that is a finite-type
  $\bar k$-algebra.  Then the structure map $\bar k \to K$ is bijective.
\end{lemma}

\begin{proof}\leanok
  By Zariski's lemma (the Nullstellensatz in Jacobson-ring form), $K$ is
  module-finite, hence integral, over $\bar k$.  An integral extension of
  an algebraically closed field is trivial, so the structure map is bijective.
\end{proof}

\begin{lemma}[Residue field at a maximal ideal of a finite-type $\bar k$-algebra]
  \label{lem:algebraMap_residueField_bijective_of_isMaximal}
  \lean{IsAlgClosed.algebraMap_residueField_bijective_of_isMaximal}
  \leanok
  \uses{lem:algebraMap_bijective_of_finiteType}
  Let $\bar k$ be algebraically closed and $A$ a finite-type $\bar k$-algebra.
  For a maximal ideal $\mathfrak p$ of $A$, the structure map
  $\bar k \to \kappa(\mathfrak p)$ into the residue field is bijective.
\end{lemma}

\begin{proof}\leanok
  The residue field $\kappa(\mathfrak{p}) = A/\mathfrak{p}$ is a surjective
  quotient of $A$, hence itself a finite-type $\bar k$-algebra; the conclusion
  follows from \cref{lem:algebraMap_bijective_of_finiteType}.
\end{proof}

\section{Residue-rationality interface}

The downstream conormal comparison reads the residue rationality not as a bare
bijection but through three formal vanishing facts ($\Omega_{K/\bar k} = 0$,
$H^1_{\mathrm{cot}}(\bar k, K) = 0$, and formal smoothness), all of which follow
once the structure map $\bar k \to K$ is bijective --- i.e.\ once $K \cong \bar k$
as a $\bar k$-algebra.  The next four blocks package exactly this interface,
keyed on bijectivity of the structure map so that the geometric application can
feed them $K = \kappa(x)$.

\begin{definition}\leanok
[$\bar k$-algebra isomorphism from a bijective structure map]
  \label{def:algEquivOfAlgebraMapBijective}
  \lean{algEquivOfAlgebraMapBijective}
  For fields $\bar k, K$ with the $\bar k$-algebra structure map $\bar k \to K$
  bijective, that map packages canonically as a $\bar k$-algebra isomorphism
  $\bar k \cong K$.
\end{definition}

\begin{lemma}\leanok
[Bijective structure map gives formal smoothness]
  \label{lem:formallySmooth_of_algebraMap_bijective}
  \lean{formallySmooth_of_algebraMap_bijective}
  \uses{def:algEquivOfAlgebraMapBijective}
  If the structure map $\bar k \to K$ is bijective, then $K$ is formally smooth
  over $\bar k$.
\end{lemma}

\begin{proof}
  \uses{def:algEquivOfAlgebraMapBijective}
  \leanok
  By \cref{def:algEquivOfAlgebraMapBijective} the structure map is a
  $\bar k$-algebra isomorphism $K \cong \bar k$.  Formal smoothness is invariant
  under algebra isomorphism, and $\bar k$ is formally smooth over itself.
\end{proof}

\begin{lemma}\leanok
[Bijective structure map kills Kähler differentials]
  \label{lem:subsingleton_kaehlerDifferential_of_algebraMap_bijective}
  \lean{subsingleton_kaehlerDifferential_of_algebraMap_bijective}
  \uses{def:algEquivOfAlgebraMapBijective}
  If the structure map $\bar k \to K$ is bijective, then $\Omega_{K/\bar k} = 0$.
\end{lemma}

\begin{proof}
  \uses{def:algEquivOfAlgebraMapBijective}
  \leanok
  Through the isomorphism $K \cong \bar k$ of
  \cref{def:algEquivOfAlgebraMapBijective}, $K$ is formally unramified over
  $\bar k$, and the module of Kähler differentials of a formally unramified
  algebra vanishes.
\end{proof}

\begin{lemma}\leanok
[Bijective structure map kills $H^1$ of the cotangent complex]
  \label{lem:subsingleton_h1Cotangent_of_algebraMap_bijective}
  \lean{subsingleton_h1Cotangent_of_algebraMap_bijective}
  \uses{lem:formallySmooth_of_algebraMap_bijective}
  If the structure map $\bar k \to K$ is bijective, then the first cotangent
  cohomology $H^1_{\mathrm{cot}}(\bar k, K)$ vanishes.
\end{lemma}

\begin{proof}
  \uses{lem:formallySmooth_of_algebraMap_bijective}
  \leanok
  $K$ is formally smooth over $\bar k$
  (\cref{lem:formallySmooth_of_algebraMap_bijective}), and the first cotangent
  cohomology of a formally smooth algebra vanishes.
\end{proof}

\section{Coheight one is closed, and the scheme-level statement}

\begin{lemma}[Coheight-one points of a dimension-$\le 1$ scheme are closed]
  \label{lem:isClosed_singleton_of_coheight_eq_one}
  \lean{AlgebraicGeometry.isClosed_singleton_of_coheight_eq_one}
  \leanok
  Let $X$ be a scheme whose underlying space has Krull dimension at most one (in
  the specialisation order).  If $x \in X$ has $\coheight(x) = 1$ then
  $\{x\}$ is closed.
\end{lemma}

\begin{proof}\leanok
  From $\krullDim X = \sup_a (\height a + \coheight a)$ and $\coheight x = 1$,
  $\dim X \le 1$ forces $\height x = 0$, i.e.\ $x$ is minimal in the
  specialisation order.  A specialisation-minimal point of a $T_0$ space is
  closed (its closure, the set of its specialisations, reduces to $\{x\}$).
\end{proof}

\begin{definition}\leanok
[$\bar k$-algebra on a residue field from the structure morphism]
  \label{def:residueFieldAlgebra}
  \lean{AlgebraicGeometry.residueFieldAlgebra}
  For a scheme $C$ over $\operatorname{Spec}\bar k$ and a point $x \in C$, the
  structure morphism $C \to \operatorname{Spec}\bar k$ induces a canonical
  $\bar k$-algebra structure on the residue field $\kappa(x)$ (via the residue-field
  map of the morphism at $x$ and the identification $\kappa(\ast)=\bar k$ at the
  base point).  This $\bar k$-algebra structure is not canonical from $\kappa(x)$
  alone but requires the data of the structure morphism, so it is recorded
  explicitly to state $\bar k$-rationality below.
\end{definition}

\begin{lemma}[Residue field of a closed point over an algebraically closed base]
  \label{lem:residueFieldIsoBase_mathlib}
  \lean{AlgebraicGeometry.residueFieldIsoBase}
  \mathlibok
  \textit{Provided by Mathlib.}
  Let $f : X \to \operatorname{Spec}\bar k$ be a morphism locally of finite type to
  the spectrum of an algebraically closed field $\bar k$, and let $x \in X$ be a
  point whose singleton $\{x\}$ is closed.  Then the residue field $\kappa(x)$ is
  canonically isomorphic, as a $\bar k$-algebra, to $\bar k$ itself.
\end{lemma}

\begin{lemma}\leanok
[Codimension-one points are $\bar k$-rational]
  \label{lem:codimOne_point_residueField_eq_kbar}
  \lean{AlgebraicGeometry.residueField_eq_of_coheight_eq_one}
  \uses{lem:isClosed_singleton_of_coheight_eq_one,
        lem:residueFieldIsoBase_mathlib,
        def:residueFieldAlgebra}
  Let $C$ be an integral curve of finite type over an algebraically closed field
  $\bar k$, and let $x \in C$ be a point with $\coheight x = 1$.  Then $x$ is a
  closed point and the canonical map $\bar k \to \kappa(x)$ is an isomorphism; in
  particular $\kappa(x) = \bar k$.
\end{lemma}

\begin{proof}
  \uses{lem:isClosed_singleton_of_coheight_eq_one,
        lem:residueFieldIsoBase_mathlib,
        def:residueFieldAlgebra}
        \leanok
  The closedness step \cref{lem:isClosed_singleton_of_coheight_eq_one} requires
  the topological dimension bound $\krullDim C \le 1$ as an explicit input,
  which the statement carries as the hypothesis $\krullDim C \le 1$.
  Granting this bound, \cref{lem:isClosed_singleton_of_coheight_eq_one} applies to
  the coheight-one point $x$ and shows that $\{x\}$ is closed; this is the first
  conclusion.

  For the residue-field rationality, the structure morphism $C \to
  \operatorname{Spec}\bar k$ is locally of finite type and $\bar k$ is
  algebraically closed, so \cref{lem:residueFieldIsoBase_mathlib} applies to the
  closed point $x$ and supplies a canonical isomorphism
  $\rho \colon \kappa(x) \cong \bar k$.  Its inverse
  $\rho^{-1} \colon \bar k \cong \kappa(x)$ is an isomorphism of rings, hence
  bijective.

  It remains to identify the canonical structure map $\bar k \to \kappa(x)$ of
  \cref{def:residueFieldAlgebra} with $\rho^{-1}$, after which bijectivity is
  immediate.  By construction the structure map is the composite of the
  residue-field map of the structure morphism at $x$ with the base identification
  $\kappa(\ast) = \bar k$.  This composite and $\rho^{-1}$ agree because they agree
  after applying the (fully faithful) spectrum functor $\operatorname{Spec}(-)$:
  there both sides reduce to the canonical morphism
  $\operatorname{Spec}\kappa(x) \to C$ from the spectrum of the residue field,
  using the compatibility of $\rho^{-1}$ with this canonical morphism and the
  compatibility of the residue-field map of the structure morphism with it.  Since
  $\operatorname{Spec}(-)$ is faithful, the two maps are equal; hence the structure
  map $\bar k \to \kappa(x)$ equals $\rho^{-1}$ and is bijective, so
  $\kappa(x) \cong \bar k$.
\end{proof}
